\documentclass[11pt]{article}

\usepackage[margin=1in]{geometry}
\usepackage{amsmath, amssymb, amsthm}
\usepackage{listings}
\usepackage{hyperref}
\usepackage{enumitem}
\usepackage{algorithm}
\usepackage{algpseudocode}
\lstset{
  language=C,
  basicstyle=\ttfamily\small,
  keywordstyle=\bfseries,
  commentstyle=\itshape,
  stringstyle=\color{blue!70!black},
  columns=fullflexible,
  keepspaces=true,
  showstringspaces=false,
  numbers=left,
  numberstyle=\scriptsize,
  frame=single,
  breaklines=true
}

\title{PETSc Data Assimilation (PetscDA) Object and the Lorenz-96 ETKF Tutorial}
\author{PETSc Documentation Draft}
\date{\today}

\begin{document}
\maketitle

\section{Overview}

PETSc's \texttt{PetscDA} (Data Assimilation) object centralizes ensemble-based data assimilation workflows. It manages ensemble storage, observation metadata, and hooks to user-provided forecast and analysis operators so that square-root ensemble filters---such as the ensemble transform Kalman filter (ETKF)---can be implemented with minimal boilerplate.

This document provides:

\begin{itemize}[nosep]
  \item A structural tour of the \texttt{PetscDA} object and its key methods.
  \item A walkthrough of the deterministic ETKF pipeline implemented in \texttt{src/da/impls/etkf/etkf.c}.
  \item An annotated explanation of the Lorenz-96 tutorial (\texttt{src/da/tutorials/ex1.c}) that exercises the PetscDA stack.
  \item Practical guidance for building, running, and extending the example.
\end{itemize}

\section{PetscDA Object Architecture}

\subsection{Lifecycle}

The \texttt{PetscDA} object follows the standard PETSc lifecycle:

\begin{enumerate}[nosep]
  \item \textbf{Creation:} \texttt{PetscDACreate(MPI\_Comm, PetscDA*)}
  \item \textbf{Configuration:} \texttt{PetscDASetType()} selects the implementation; \texttt{PETSCDAETKF} is the built-in deterministic ETKF backend.
  \item \textbf{Sizing:} \texttt{PetscDASetSizes()} records the state dimension, observation dimension, and ensemble size (all logically collective).
  \item \textbf{Assembly:} \texttt{PetscDASetUp()} allocates the dense ensemble matrix, observation-error vector, and orthogonal transform workspace.
  \item \textbf{Usage:} Ensemble members are accessed via \texttt{PetscDAGetEnsembleMember()}, \texttt{PetscDASetEnsembleMember()}, and analyzed via \texttt{PetscDAComputeMean()}, \texttt{PetscDAComputeAnomalies()}, \texttt{PetscDAAnalysis()}, and \texttt{PetscDAApplyModel()}.
  \item \textbf{Destruction:} \texttt{PetscDADestroy()} releases all internal allocations.
\end{enumerate}

The \texttt{PetscDA} virtual table (\texttt{struct \_PetscDAOps}) carries function pointers for analysis, model application, mean/anomaly computation, option parsing, and destruction. Implementations provide these hooks when they register with \texttt{PetscDARegister()}.

\subsection{Ensemble Storage}

Internally, \texttt{PetscDA} maintains:

\begin{itemize}[nosep]
  \item \texttt{Mat ensemble}: a dense matrix of size \texttt{state\_size} $\times$ \texttt{ensemble\_size}. Column $j$ stores ensemble member $x^{(j)}$.
  \item \texttt{Vec obs\_error\_var}: observation-error variances (for diagonal $R$ models).
  \item \texttt{Mat U}: an ensemble-space orthogonal transform (defaults to identity).
\end{itemize}

Helpers such as \texttt{PetscDAComputeMean()} and \texttt{PetscDAComputeAnomalies()} return PETSc vectors/matrices borrowed from these internal structures, preserving reference counts and allowing zero-copy manipulation.

\section{ETKF Implementation Sketch}

\subsection{Analysis Flow (Algorithm 6.4)}

The deterministic ETKF analysis in \texttt{PetscDAETKFAnalysis()} proceeds as follows:


\begin{algorithm}
\caption{Ensemble Transform Kalman Filter (ETKF) - Deterministic}
\begin{algorithmic}[1]
\State \textbf{Initialize:} Ensemble $\mathbf{E}_0 = [\mathbf{x}_0^{(1)}, \ldots, \mathbf{x}_0^{(m)}]$ \Comment{$\mathbf{E}_0 \in \mathbb{R}^{n \times m}$}
\For{$k = 1, 2, \ldots, K$}
    \State \textbf{// Forecast Step}
    \For{$i = 1$ to $m$}
        \State $\mathbf{x}_k^{f,(i)} = \mathcal{M}_{k-1}(\mathbf{x}_{k-1}^{a,(i)})$ \Comment{$\mathbf{x}_k^{f,(i)} \in \mathbb{R}^{n}$}
    \EndFor
    \State Assemble forecast matrix: $\mathbf{E}_k^f = [\mathbf{x}_k^{f,(1)}, \ldots, \mathbf{x}_k^{f,(m)}]$ \Comment{$\mathbf{E}_k^f \in \mathbb{R}^{n \times m}$}
    \State
    \State \textbf{// Analysis Step (if observation available)}
    \If{observation $\mathbf{y}_k$ available}
        \State Compute ensemble mean: $\overline{\mathbf{x}}_k^f = \frac{1}{m}\sum_{i=1}^m \mathbf{x}_k^{f,(i)}$ \Comment{$\overline{\mathbf{x}}_k^f \in \mathbb{R}^{n}$}
        \State Compute anomalies: $\mathbf{X}_k = \frac{1}{\sqrt{m-1}}(\mathbf{E}_k^f - \overline{\mathbf{x}}_k^f \mathbf{1}^T)$ \Comment{$\mathbf{X}_k \in \mathbb{R}^{n \times m}$}
        \State Apply observation operator: $\mathbf{Z}_k = \mathcal{H}(\mathbf{E}_k^f)$ \Comment{$\mathbf{Z}_k \in \mathbb{R}^{b \times m}$}
        \State Compute obs ensemble mean: $\overline{\mathbf{y}}_k = \frac{1}{m}\sum_{i=1}^m \mathcal{H}(\mathbf{x}_k^{f,(i)})$ \Comment{$\overline{\mathbf{y}}_k \in \mathbb{R}^{b}$}
        \State Compute obs anomalies: $\mathbf{S}_k = \frac{1}{\sqrt{m-1}}\mathbf{R}^{-1/2}(\mathbf{Z}_k - \overline{\mathbf{y}}_k\mathbf{1}^T)$ \Comment{$\mathbf{S}_k \in \mathbb{R}^{b \times m}$}
        \State Compute raw innovation: $\boldsymbol{\delta}_k = \mathbf{y}_k - \overline{\mathbf{y}}_k$ \Comment{$\boldsymbol{\delta}_k \in \mathbb{R}^{b}$}
        \State Whiten innovation: $\tilde{\boldsymbol{\delta}}_k = \mathbf{R}^{-1/2}\boldsymbol{\delta}_k$ \Comment{$\tilde{\boldsymbol{\delta}}_k \in \mathbb{R}^{b}$}
        \State Compute transform matrix: $\mathbf{T}_k = (\mathbf{I}_m + \mathbf{S}_k^T\mathbf{S}_k)^{-1}$ \Comment{$\mathbf{T}_k \in \mathbb{R}^{m \times m}$}
        \State Compute weight vector: $\mathbf{w}_k = \mathbf{T}_k \mathbf{S}_k^T \tilde{\boldsymbol{\delta}}_k$ \Comment{$\mathbf{w}_k \in \mathbb{R}^{m}$}
        \State Form deterministic map: $\mathbf{G}_k = \mathbf{w}_k\mathbf{1}^T + \sqrt{m-1}\,\mathbf{T}_k^{1/2}\mathbf{U}$ \Comment{$\mathbf{G}_k \in \mathbb{R}^{m \times m}$}
        \State Update ensemble: $\mathbf{E}_k^a = \overline{\mathbf{x}}_k^f\mathbf{1}^T + \mathbf{X}_k \mathbf{G}_k$ \Comment{$\mathbf{E}_k^a \in \mathbb{R}^{n \times m}$}
        \State where $\mathbf{U}$ is orthogonal with $\mathbf{U}\mathbf{1} = \mathbf{1}$ (the PETSc implementation stores it as \texttt{da->U} and defaults to $\mathbf{I}_m$)
    \Else
        \State $\mathbf{E}_k^a = \mathbf{E}_k^f$ \Comment{$\mathbf{E}_k^a \in \mathbb{R}^{n \times m}$}
    \EndIf
\EndFor
\end{algorithmic}
\end{algorithm}

Throughout these steps, \texttt{PetscUseTypeMethod()} dispatches to the configured backend, enforcing type-safety and consistent error propagation.

\subsection{Model Propagation}

\texttt{PetscDAApplyModel()} iterates over ensemble members, calling a user-supplied function $M(x, x_{\text{new}}; \text{ctx})$ that advances one state. The Lorenz-96 tutorial uses an explicit Runge--Kutta 4 (\texttt{TSRK4}) step to propagate each member.

\section{Lorenz-96 Tutorial (\texttt{ex1.c})}

\subsection{Problem Setup}

\begin{itemize}[nosep]
  \item \textbf{Model:} One-dimensional Lorenz-96 system with forcing $F=8$, state dimension $n=40$, periodic boundary conditions.
  \item \textbf{Truth State:} Generated by integrating the same model with identical forcing.
  \item \textbf{Observations:} Identity operator with additive Gaussian noise of standard deviation \texttt{-obs\_error}.
  \item \textbf{Ensemble:} Initialized around the climatological state (all entries $=F$) with Gaussian perturbations.
  \item \textbf{Assimilation:} Deterministic ETKF using either Cholesky or eigen-based square roots (\texttt{-petscdaetkf\_sqrt\_type}).
\end{itemize}

\subsection{Key Code Fragments}

\paragraph{Lorenz-96 Tendency.}
\begin{lstlisting}
for (i = xs; i < xs + xm; i++) {
  f[i] = (x[i + 1] - x[i - 2]) * x[i - 1] - x[i] + l95->F;
}
\end{lstlisting}
The DMDA ghost region supplies the necessary cyclic neighbors.

\paragraph{Single-Step Integrator.}
\begin{lstlisting}
PetscCall(TSCreate(PETSC_COMM_SELF, &ts));
PetscCall(TSSetType(ts, TSRK));
PetscCall(TSRKSetType(ts, TSRK4));
PetscCall(TSSetTimeStep(ts, l95->dt));
PetscCall(TSSetMaxSteps(ts, 1));
PetscCall(TSSolve(ts, x_out));
\end{lstlisting}

\paragraph{Ensemble Initialization.}
\begin{lstlisting}
for (PetscInt i = 0; i < ensemble_size; i++) {
  VecCopy(x0, member);
  VecSetRandomGaussian(spread, rng, 0.0, obs_error_std);
  VecAXPY(member, 1.0, spread);
  PetscDASetEnsembleMember(da_ctx, i, member);
}
\end{lstlisting}

\paragraph{Forecast/Analysis Loop.}
\begin{lstlisting}
for (step = 0; step <= steps; step++) {
  PetscDAComputeMean(da_ctx, x_mean);
  VecCopy(x_mean, x_forecast);
  /* Compute RMSE against truth ... */

  if (step % obs_freq == 0) {
    VecSetRandomGaussian(obs_noise, rng, 0.0, obs_error_std);
    VecCopy(truth_state, observation);
    VecAXPY(observation, 1.0, obs_noise);
    PetscDAAnalysis(da_ctx, observation, Lorenz96ObsIdentity, NULL);
  }

  if (step < steps) {
    PetscDAApplyModel(da_ctx, Lorenz96Step, &l95_ctx);
    Lorenz96Step(truth_state, truth_next, &l95_ctx);
    VecCopy(truth_next, truth_state);
  }
}
\end{lstlisting}

\subsection{Diagnostics}

\texttt{ex1.c} reports:

\begin{itemize}[nosep]
  \item Instantaneous forecast and analysis root-mean-square error (RMSE) against the truth state.
  \item Post burn-in averages of forecast and analysis RMSE.
  \item Number of assimilated observations.
\end{itemize}

Users can enable matrix/vector views with \texttt{-petscda\_view} and toggle square-root methods via \texttt{-petscdaetkf\_sqrt\_type cholesky|eigen}.

\subsection{Running the Example}

\begin{lstlisting}[language=bash]
$ petsc/bin/petsc_tsexamples -ts_tutorial_ex1 \
    -steps 120 -burn 10 -obs_freq 2 -obs_error 0.5 \
    -ensemble_size 40 -petscda_view
\end{lstlisting}

The options can be tuned to explore ensemble size sensitivity, observation frequency, or alternative square-root factorizations.

\section{Extending the Tutorial}

\paragraph{Alternative Models.} Replace \texttt{Lorenz96Step} with a domain-specific model. The model callback must obey the signature \texttt{PetscErrorCode (*model)(Vec, Vec, void*)}.

\paragraph{Custom Observation Operators.} Implement \texttt{observation\_operator(Vec state, Vec prediction, void *ctx)} to map states into observation space; update \texttt{PetscDAAnalysis()} accordingly.

\paragraph{Non-diagonal Observation Errors.} Extend the PetscDA backend to accept general covariance matrices by replacing \texttt{VecPointwiseMult} operations with solves against $R^{1/2}$.

\paragraph{Monitoring.} Attach custom logging via \texttt{PetscDASetOptionsPrefix()} and PETSc's \texttt{-log\_view} framework to profile matrix factorizations and communication costs.

\section{Summary}

The \texttt{PetscDA} object abstracts the repetitive scaffolding of ensemble-based assimilation algorithms. The Lorenz-96 ETKF tutorial demonstrates the end-to-end workflow:

\begin{enumerate}[nosep]
  \item Configure ensemble sizes and model parameters.
  \item Initialize the truth trajectory and ensemble.
  \item Run a forecast/analysis loop using \texttt{PetscDAApplyModel()} and \texttt{PetscDAAnalysis()}.
  \item Evaluate performance via RMSE diagnostics.
\end{enumerate}

From here, users can adapt the example to realistic models, incorporate domain-specific observation operators, and experiment with advanced square-root filters---all while reusing PETSc's scalable linear algebra and time-integration components.

\end{document}